% Part: first-order-logic
% Chapter: axiomatic-deduction
% Section: axioms-rules-propositional

\documentclass[../../../include/open-logic-section]{subfiles}

\begin{document}

\iftag{FOL}
      {\olfileid{fol}{axd}{prp}}
      {\olfileid{pl}{axd}{prp}}

\olsection{विधानीय संयोजकांसाठी स्वयंसिद्धके आणि नियम}

\begin{defn}[स्वयंसिद्धके]
विधानीय संयोजकांसाठीच्या \emph{स्वयंसिद्धकांचा} संच $\PAx$ हा पुढील रूपांच्या
सर्व !!{formula}s चा बनलेला आहे:
\begin{align}
  & (!A \land !B) \lif !A \ollabel{ax:land1}\\
  & (!A \land !B) \lif !B \ollabel{ax:land2}\\
  & !A \lif (!B \lif (!A \land !B)) \ollabel{ax:land3}\\
  & !A \lif (!A \lor !B) \ollabel{ax:lor1}\\
  & !A \lif (!B \lor !A) \ollabel{ax:lor2}\\
  & (!A \lif !C) \lif ((!B \lif !C) \lif ((!A \lor !B) \lif !C)) \ollabel{ax:lor3}\\
  & !A \lif (!B \lif !A) \ollabel{ax:lif1}\\
  & (!A \lif (!B \lif !C)) \lif ((!A \lif !B) \lif (!A \lif !C)) \ollabel{ax:lif2}\\
  & (!A \lif !B) \lif ((!A \lif \lnot !B) \lif \lnot !A) \ollabel{ax:lnot1}\\
  & \lnot !A \lif (!A \lif !B) \ollabel{ax:lnot2}\\
  & \ltrue \ollabel{ax:ltrue}\\
  & \lfalse \lif !A \ollabel{ax:lfalse1}\\
  & (!A \lif \lfalse) \lif \lnot !A \ollabel{ax:lfalse2}\\
  & \lnot\lnot !A \lif !A \ollabel{ax:dne}
\end{align}
\end{defn}

\begin{defn}[विध्यात्मक अनुमान]
  जर $!B$ आणि $!B \lif !A$ हे !!a{derivation} मध्ये आधीच आलेले असतील,
  तर $!A$ ही निगमनाची योग्य पायरी असते.
\end{defn}

विध्यात्मक अनुमान या नियमाचा संक्षेप आपण ``\MP'' असा करू.

\end{document}
